% Part: first-order-logic
% Chapter: beyond
% Section: second-order-logic

\documentclass[../../../include/open-logic-section]{subfiles}

\begin{document}

\olfileid{fol}{byd}{sol}

\olsection{ద్వితీయ-స్థాయి తర్కం}

ద్వితీయ-స్థాయి తర్కపు భాషలో వ్యక్తుల \tetoken{ఒక వ్యక్తి క్షేత్రంపై}{!!a{domain}} మాత్రమే కాకుండా, ఆ
\tetoken{వ్యక్తి క్షేత్రంపై}{!!{domain}} సంబంధాలపైనా పరిమాణీకరించవచ్చు. ఒక మొదటిస్థాయి
భాష~$\Lang{L}$ ఇచ్చినప్పుడు, ప్రతి $k$కు $k$-స్థాన సంబంధాలపై విస్తరించే
\tetoken{చరాలు}{!!{variable}s} $R$ను చేర్చి, ఆ చరాలపై పరిమాణీకరణను అనుమతిస్తాం. $R$ ఒక
$k$-స్థాన సంబంధపు \tetoken{ఒక చరం}{!!a{variable}}, $t_1$, \dots,~$t_k$లు సాధారణ
(మొదటిస్థాయి) పదాలు అయితే, $\Atom{R}{t_1,\dots,t_k}$ ఒక పరమాణు
\tetoken{సూత్రం}{!!{formula}}. మిగతా విషయాల్లో, ద్వితీయ-స్థాయి పరిమాణీకరణకు అదనపు
నిబంధనలు చేర్చి, మొదటిస్థాయి తర్కంలోలాగే \tetoken{సూత్రాలు}{!!{formula}s} సమితిని
నిర్వచిస్తాం. మనకు మొదటిస్థాయి పదాలకు మాత్రమే \tetoken{తాదాత్మ్య విధేయం}{!!{identity}} ఉందని
గమనించండి: $R$, $S$లు ఒకే అరిటీ~$k$ గల సంబంధ \tetoken{చరాలు}{!!{variable}s} అయితే,
$\eq[R][S]$ను కిందిదానికి సంక్షిప్త రూపంగా నిర్వచించవచ్చు:
\[
\lforall[x_1 \dots][\lforall[x_k][(\Atom{R}{x_1, \dots, x_k} \liff
  \Atom{S}{x_1, \dots, x_k})]].
\]

ద్వితీయ-స్థాయి తర్కపు నియమాలు పరిమాణక నియమాలను కొత్త ద్వితీయ-స్థాయి
చరాలకు విస్తరిస్తాయి. అయితే ఇక్కడ, ఈ చరాలు $\Lang{L}$లోని
\tetoken{విధేయ సంకేతాలతో}{!!{predicate}s}, మరింత సాధారణంగా $\Lang{L}$లోని \tetoken{సూత్రాలతో}{!!{formula}s} ఎలా
పరస్పరచర్య చేస్తాయో కొంచెం జాగ్రత్తగా వివరించాలి. కనీసం, సంబంధ చరాలను
పదాలుగా పరిగణిస్తాం; అందువల్ల కింది రూపంలోని నిగమనాలు ఉంటాయి:
\[
!A(R) \Proves \lexists[R][!A(R)]
\]
కానీ $\Lang{L}$, స్థిర సంబంధ సంకేతం~$<$ గల అంకగణిత భాష అయితే, కింది
నిగమనం కూడా చెల్లుబాటవుతుందని ఆశిస్తాం:
\[
x < y \Proves \lexists[R][\Atom{R}{x,y}]
\]
లేదా ఇచ్చిన \tetoken{సూత్రం}{!!{formula}}~$!A$కు,
\[
!A(x_1, \dots, x_k) \Proves \lexists[R][\Atom{R}{x_1,\dots,x_k}]
\]
మరింత సాధారణంగా, కింది రూపంలోని నిగమనాలను అనుమతించాలనుకోవచ్చు:
\[
\Subst{!A}{\lambd[\vec x][!B(\vec x)]}{R} \Proves
\lexists[R][!A]
\]
ఇక్కడ $\Subst{!A}{\lambd[\vec x][!B(\vec x)]}{R}$ అనేది, $!A$లో
$\Atom{R}{t_1,\dots,t_k}$ రూపంలోని ప్రతి పరమాణు \tetoken{సూత్రాన్ని}{!!{formula}}
$!B(t_1, \dots, t_k)$తో ప్రతిస్థాపించడం వల్ల వచ్చే ఫలితాన్ని
సూచిస్తుంది. ఈ చివరి నియమం ఒక {\em ధర్మసంగ్రహ పథకం}, అంటే కింది
రూపంలోని స్వీకృతం ఉండటానికి సమానం:
\[
\lexists[R][\lforall[x_1, \dots, x_k][(!A(x_1, \dots, x_k) \liff
\Atom{R}{x_1, \dots, x_k})]],
\]
ద్వితీయ-స్థాయి భాషలో $R$ స్వేచ్ఛా చరం కాని ప్రతి \tetoken{సూత్రం}{!!{formula}}~$!A$కు
ఇటువంటి ఒక స్వీకృతం ఉంటుంది. (అభ్యాసం: $!A$లో $R$ను సంభవించనిస్తే ఈ
పథకం వైరుధ్యభరితమని చూపండి!)
\sourcecorrection{OLTEFOLBYD-002}{ప్రతిస్థాపించాల్సిన పరమాణు సూత్రాన్ని ఆంగ్ల మూలం \texttt{\string\Obj\{R\}\{t_1,\dots,t_k\}}గా ముద్రించింది; అదే పేరాలోని పరమాణు-సూత్ర నిబంధనకు అనుగుణంగా \texttt{\string\Atom\{R\}\{t_1,\dots,t_k\}}గా సరిచేశాం.}

తర్కశాస్త్రవేత్తలు ``ద్వితీయ-స్థాయి తర్కపు స్వీకృతాలు'' అని
చెప్పినప్పుడు, సాధారణంగా ధర్మసంగ్రహ పథకంతోపాటు ద్వితీయ-స్థాయి
పరిమాణక నియమాల ద్వారా మొదటిస్థాయి తర్కానికి చేసిన కనిష్ఠ విస్తరణను
ఉద్దేశిస్తారు. అయితే ఈ స్వీకృతాలు, నియమాల బలహీన ఉపవ్యవస్థలను
అధ్యయనం చేయడం తరచుగా ఆసక్తికరం. ఉదాహరణకు, పూర్తి సాధారణతలో
ధర్మసంగ్రహ స్వీకృత పథకం \emph{ఇంప్రెడికేటివ్}: ద్వితీయ-స్థాయి
పరిమాణకాలున్న \tetoken{ఒక సూత్రం}{!!a{formula}} ద్వారా ``నిర్వచితమైన''
$\Atom{R}{x_1, \dots, x_k}$ అనే సంబంధం ఉందని ప్రకటించడానికి ఇది
అనుమతిస్తుంది; ఈ పరిమాణకాలు అటువంటి సంబంధాలన్నింటి సమితిపై
విస్తరిస్తాయి---ఆ సమితిలో $R$ కూడా ఉంటుంది!{} ఇరవయ్యవ శతాబ్దపు
ఆరంభం నాటికి, రసెల్ వైరుధ్యానికి ఇటువంటి నిర్వచనాలే కారణమని భావించి,
గణిత పునాదులను అభివృద్ధి చేసేటప్పుడు వాటిని నివారించడం ఒక సాధారణ
ప్రతిస్పందనగా ఉండేది. \tetoken{సూత్రం}{!!{formula}}~$!A$లో ద్వితీయ-స్థాయి పరిమాణకాలను
వాడకుండా నిషేధిస్తే, కొంత బలహీనమైన {\em ప్రెడికేటివ్} ధర్మసంగ్రహ
రూపం లభిస్తుంది.

అర్థపర దృక్కోణం నుంచి, భాషకు ఒక మొదటిస్థాయి \tetoken{నిర్మాణంతో}{!!{structure}}పాటు,
ద్వితీయ-స్థాయి పరిమాణకాలు విస్తరించే ఆ \tetoken{వ్యక్తి క్షేత్రంపై}{!!{domain}} సంబంధాల సమితి
కలసి ద్వితీయ-స్థాయి \tetoken{నిర్మాణం}{!!{structure}} ఏర్పడుతుందని భావించవచ్చు (మరింత
ఖచ్చితంగా, ప్రతి $k$కు అరిటీ~$k$ గల సంబంధాల సమితి ఒకటి ఉంటుంది).
ధర్మసంగ్రహాన్ని \tetoken{వ్యుత్పత్తి}{!!{derivation}} వ్యవస్థలో చేర్చితే, ధర్మసంగ్రహ
స్వీకృతాలను సంతృప్తిపరచడానికి ``ద్వితీయ-స్థాయి భాగంలో'' తగినన్ని
సంబంధాలు ఉండాలనే అదనపు నిబంధన వస్తుంది---లేకపోతే \tetoken{వ్యుత్పత్తి}{!!{derivation}}
వ్యవస్థ యథార్థం కాదు!{} తగినన్ని సంబంధాలు ఉన్నాయని నిర్ధారించడానికి
ఒక సులభమైన మార్గం, మొదటిస్థాయి భాగంపై ఉన్న సంబంధాలన్నింటినీ
ద్వితీయ-స్థాయి భాగంగా తీసుకోవడం. అటువంటి \tetoken{ఒక నిర్మాణాన్ని}{!!a{structure}}
\emph{పూర్ణం} అంటారు; ఒక అర్థంలో భాషకు నిజమైన ``ఉద్దేశిత
\tetoken{నిర్మాణం}{!!{structure}}'' అదే. మన దృష్టిని పూర్ణ \tetoken{నిర్మాణాలు}{!!{structure}s}కే పరిమితం
చేస్తే, దానిని \emph{పూర్ణ} ద్వితీయ-స్థాయి అర్థవిచారం అంటారు. అప్పుడు
\tetoken{నిర్మాణాన్ని}{!!{structure}} నిర్దేశించడం అంటే మొదటిస్థాయి భాగాన్ని నిర్దేశించడమే;
దాని నుంచి ద్వితీయ-స్థాయి భాగపు విషయం పరోక్షంగా నిర్ణయమవుతుంది.

సంగ్రహంగా, ద్వితీయ-స్థాయి తర్కం గురించి మాట్లాడేటప్పుడు కొంత
సందిగ్ధత ఉంటుంది. \tetoken{వ్యుత్పత్తి}{!!{derivation}} వ్యవస్థ విషయంలో కింది రెండింటిలో
ఏదో ఒకదాన్ని ఉద్దేశించవచ్చు:
\begin{enumerate}
\item కొన్ని ధర్మసంగ్రహ స్వీకృతాలతో కూడిన ``కనిష్ఠ'' ద్వితీయ-స్థాయి
  \tetoken{వ్యుత్పత్తి}{!!{derivation}} వ్యవస్థ.
\item పూర్ణ ధర్మసంగ్రహంతో కూడిన ``ప్రామాణిక'' ద్వితీయ-స్థాయి
  \tetoken{వ్యుత్పత్తి}{!!{derivation}} వ్యవస్థ.
\end{enumerate}
అర్థవిచారం విషయంలో కింది రెండింటిలో ఏదో ఒకదానిపై ఆసక్తి ఉండవచ్చు:
\begin{enumerate}
\item ఒక మొదటిస్థాయి భాగం, ధర్మసంగ్రహ స్వీకృతాలను సంతృప్తిపరచగలంత
  పెద్ద ద్వితీయ-స్థాయి భాగం కలసి \tetoken{ఒక నిర్మాణం}{!!a{structure}} అయ్యే ``బలహీన''
  అర్థవిచారం.
\item పూర్ణ \tetoken{నిర్మాణాలను}{!!{structure}s} మాత్రమే పరిగణించే ``ప్రామాణిక''
  ద్వితీయ-స్థాయి అర్థవిచారం.
\end{enumerate}
తర్కశాస్త్రవేత్తలు తాము ఉద్దేశించే \tetoken{వ్యుత్పత్తి}{!!{derivation}} వ్యవస్థను లేదా
అర్థవిచారాన్ని ప్రత్యేకంగా చెప్పనప్పుడు, సాధారణంగా ప్రతి జాబితాలోని
రెండవ అంశాన్నే సూచిస్తారు. ఈ అర్థవిచారం వల్ల కలిగే ప్రయోజనం ఏమిటంటే,
మనం చూడబోయే విధంగా అనేక సహజ గణిత నిర్మాణాలకు ఇది సమరూపత వరకు
ఏకైకమైన (కేటగారికల్) వర్ణనలను ఇస్తుంది; అదే సమయంలో \tetoken{వ్యుత్పత్తి}{!!{derivation}}
వ్యవస్థ చాలా బలమైనది, ఈ అర్థవిచారానికి యథార్థమైనది. లోపం ఏమిటంటే,
\tetoken{వ్యుత్పత్తి}{!!{derivation}} వ్యవస్థ ఈ అర్థవిచారానికి \emph{సంపూర్ణమైనది కాదు};
నిజానికి ప్రభావవంతంగా ఇచ్చిన \emph{ఏ} \tetoken{వ్యుత్పత్తి}{!!{derivation}} వ్యవస్థా పూర్ణ
ద్వితీయ-స్థాయి అర్థవిచారానికి సంపూర్ణం కాదు. మరోవైపు, బలహీన
అర్థవిచారానికి \tetoken{వ్యుత్పత్తి}{!!{derivation}} వ్యవస్థ \emph{సంపూర్ణమే} అని చూస్తాం;
దీనివల్ల ఒక వాక్యం నిరూపించదగినది కాకపోతే, అది అసత్యమయ్యే
\emph{ఏదో ఒక} \tetoken{నిర్మాణం}{!!{structure}}---అది పూర్ణ నిర్మాణం కావలసిన అవసరం
లేదు---ఉంటుందని తెలుస్తుంది.

ద్వితీయ-స్థాయి తర్కపు భాష చాలా సమృద్ధమైనది. ఏకస్థాన సంబంధాలను
\tetoken{వ్యక్తి క్షేత్రం}{!!{domain}} యొక్క ఉపసమితులతో గుర్తించవచ్చు; అందువల్ల ముఖ్యంగా ఈ
సమితులపై పరిమాణీకరించవచ్చు. ఉదాహరణకు, సహజ సంఖ్యల ఆగమనాన్ని ఒకే
స్వీకృతంతో వ్యక్తం చేయవచ్చు:
\[
\lforall[R][((\Atom{R}{\Obj{0}} \land \lforall[x][(\Atom{R}{x} \lif
    \Atom{R}{x'})]) \lif \lforall[x][\Atom{R}{x}])].
\]
అంకగణిత భాషలో $\Obj 0, \Obj \prime, +, \times$, $<$ సంకేతాలు
ఉన్నాయని తీసుకుంటే, వాటి ప్రవర్తనను వివరించడానికి కింది స్వీకృతాలను
చేర్చవచ్చు:
\begin{enumerate}
\item $\lforall[x][\lnot x' = \Obj 0]$
\item $\lforall[x][\lforall[y][(x' = y' \lif x = y)]]$
\item $\lforall[x][(x + \Obj 0) = x]$
\item $\lforall[x][\lforall[y][(x + y') = (x + y)']]$
\item $\lforall[x][(x \times \Obj 0) = \Obj 0]$
\item $\lforall[x][\lforall[y][(x \times y') = ((x \times y) + x)]]$
\item $\lforall[x][\lforall[y][(x < y \liff \lexists[z][y = (x + z')])]]$
\end{enumerate}
\sourcecorrection{OLTEFOLBYD-003}{భాషకు ఇచ్చిన ఉత్తరాధికారి సంకేతం $\Obj\prime$ అయినప్పటికీ రెండవ స్వీకృతంలో ఆంగ్ల మూలం నిర్వచించని $s$ను వాడింది; పక్కనున్న స్వీకృతాలు, తరువాతి సమరూపత నిరూపణకు అనుగుణంగా దానిని $x'=y'$గా సరిచేశాం.}
మనం పూర్ణ ద్వితీయ-స్థాయి అర్థవిచారాన్ని వాడుతున్నామనే షరతుతో, ఈ
స్వీకృతాలు పైనున్న ఆగమన స్వీకృతంతో కలసి అంకగణితపు ప్రామాణిక
నమూనా~$\Struct{N}$ అనే \tetoken{నిర్మాణానికి}{!!{structure}} సమరూపత వరకు ఏకైకమైన వర్ణనను
ఇస్తాయని చూపడం కష్టం కాదు. ఈ స్వీకృతాలు సత్యమయ్యే ఏ
\tetoken{నిర్మాణం}{!!{structure}}~$\Struct{M}$ ఇచ్చినా, $\Nat$పై సాధారణ పునరావర్తనాన్ని
వాడి $\Nat$ నుంచి $\Struct{M}$ యొక్క \tetoken{వ్యక్తి క్షేత్రానికి}{!!{domain}} $f$ అనే ప్రమేయాన్ని
నిర్వచించండి; ఇక్కడ $f(0) = \Assign{\Obj 0}{M}$,
$f(x+1) = \Assign{\prime}{M}(f(x))$. $\Nat$పై సాధారణ ఆగమనాన్ని,
$\Struct M$లో స్వీకృతాలు (1),~(2) సత్యమనే విషయాన్ని వాడితే $f$
\tetoken{ఒకటి-ఒకటి}{!!{injective}} అని తెలుస్తుంది. $f$ \tetoken{సంగ్రస్త}{!!{surjective}} అని చూడటానికి,
$f$ వ్యాప్తిలో ఉన్న $\Domain{M}$ మూలకాల సమితిని~$P$గా తీసుకోండి.
$\Struct M$ పూర్ణమైనది కాబట్టి $P$ ద్వితీయ-స్థాయి \tetoken{వ్యక్తి క్షేత్రంలో}{!!{domain}}
ఉంటుంది. $f$ నిర్మాణం వల్ల $\Assign{\Obj 0}{M}$ అనేది~$P$లో
ఉంటుందని, $\Assign{\prime}{M}$ కింద $P$ సంవృతమని తెలుసు.
$\Struct M$లో ఆగమన స్వీకృతం (ముఖ్యంగా~$P$కు) సత్యమనే విషయం, $P$
అనేది $\Struct M$ యొక్క మొత్తం మొదటిస్థాయి \tetoken{వ్యక్తి క్షేత్రానికి}{!!{domain}} సమానమని
హామీ ఇస్తుంది. దీనివల్ల $f$ \tetoken{ఒక ద్విగుణ ప్రమేయం}{!!a{bijection}} అని తెలుస్తుంది. $\Nat$పై
సాధారణ ఆగమనాన్ని పదేపదే వాడి, $f$ ఒక నిర్మాణ-పరిరక్షక సంవిధానం
(హోమోమార్ఫిజం) అని చూపడం కూడా ఇంతకంటే కష్టం కాదు.

సమితి-సిద్ధాంత పరంగా ప్రమేయం ఒక ప్రత్యేక రకమైన సంబంధం మాత్రమే;
ఉదాహరణకు, ఏకస్థాన ప్రమేయం~$f$ను
$\lforall[x][\lexists![y][R(x,y)]]$ను సంతృప్తిపరచే ద్విస్థాన
సంబంధం~$R$తో గుర్తించవచ్చు. ఫలితంగా ప్రమేయాలపైనా పరిమాణీకరించవచ్చు.
పూర్ణ అర్థవిచారాన్ని వాడి, $\Struct M$ యొక్క \tetoken{వ్యక్తి క్షేత్రం}{!!{domain}} నుంచి దాని యథార్థ
ఉపసమితిలోకి ఒకటి-ఒకటి ప్రమేయం గల \tetoken{నిర్మాణాలు}{!!{structure}s} $\Struct M$ యొక్క
తరగతినే అనంత \tetoken{నిర్మాణాలు}{!!{structure}s} తరగతిగా నిర్వచించవచ్చు:
\[
\lexists[f][(\lforall[x][\lforall[y][(\eq[f(x)][f(y)] \lif \eq[x][y])]]
  \land \lexists[y][\lforall[x][\eq/[f(x)][y]]])].
\]
ఈ వాక్యపు నిషేధం అప్పుడు పరిమిత \tetoken{నిర్మాణాలు}{!!{structure}s} తరగతిని నిర్వచిస్తుంది.

అంతేకాక, రేఖీయ క్రమ నిర్వచనానికి కిందిదాన్ని చేర్చి సుక్రమాల
తరగతిని నిర్వచించవచ్చు:
\[
\lforall[P][(\lexists[x][\Atom{P}{x}] \lif \lexists[x][(\Atom{P}{x}
    \land \lforall[y][(y < x \lif \lnot \Atom{P}{y})])])].
\]
``సమితి''ని ``ఏకస్థాన సంబంధం''తో గుర్తించడాన్ని మినహాయిస్తే, ఖాళీ
కాని ప్రతి సమితికీ ఒక కనిష్ఠ మూలకం ఉందని ఇది ప్రకటిస్తుంది. మరో
ఉదాహరణగా, శీర్షాలను పరస్పరం కలవని భాగాలుగా తుచ్ఛం కాని విధంగా
విభజించడం సాధ్యం కాదని చెప్పడం ద్వారా గ్రాఫ్‌ల సంధానితత్వ భావనను
వ్యక్తం చేయవచ్చు:
\[
\lnot \lexists[A][(\lexists[x][A(x)] \land \lexists[y][\lnot A(y)]
  \land \lforall[w][\lforall[z][((\Atom{A}{w} \land \lnot \Atom{A}{z})
      \lif \lnot \Atom{R}{w,z})]])].
\]
మరొక ఉదాహరణగా, \tetoken{వ్యక్తి క్షేత్రం}{!!{domain}} పరిమాణం సరిసంఖ్యగా ఉండే పరిమిత
\tetoken{నిర్మాణాలు}{!!{structure}s} తరగతిని ఒక అభ్యాసంగా నిర్వచించడానికి ప్రయత్నించవచ్చు.
ఇంకా విశేషంగా, కరణీయ సంఖ్యలను కలిగిన సంపూర్ణ క్రమిత క్షేత్రంగా
వాస్తవ సంఖ్యలకు సమరూపత వరకు ఏకైకమైన వర్ణన ఇవ్వవచ్చు.

సంక్షిప్తంగా, మొదటిస్థాయి తర్కంకంటే ద్వితీయ-స్థాయి తర్కం ఎంతో ఎక్కువ
వ్యక్తీకరణశక్తి గలది. ఇది శుభవార్త; ఇక చెడు వార్త. పూర్ణ
ద్వితీయ-స్థాయి అర్థవిచారానికి సంపూర్ణమైన ప్రభావవంతమైన
\tetoken{వ్యుత్పత్తి}{!!{derivation}} వ్యవస్థ ఏదీ లేదని ఇప్పటికే పేర్కొన్నాం. మంచికో చెడుకో,
సంగ్రహత్వం, లొవెన్‌హైమ్--స్కోలెమ్ సిద్ధాంతాలతో సహా మొదటిస్థాయి
తర్కపు అనేక ధర్మాలు ఇక్కడ లేవు.

మరోవైపు, పూర్ణ ద్వితీయ-స్థాయి అర్థవిచారాన్ని వదిలి బలహీన
అర్థవిచారాన్ని తీసుకోవడానికి సిద్ధమైతే, కనిష్ఠ ద్వితీయ-స్థాయి
\tetoken{వ్యుత్పత్తి}{!!{derivation}} వ్యవస్థ ఈ అర్థవిచారానికి సంపూర్ణం. మరో మాటలో,
$\Proves$ను ``కనిష్ఠ వ్యవస్థలో నిరూపిస్తుంది'' అని, $\Entails$ను
``బలహీన అర్థవిచారంలో తార్కికంగా అనుగమిస్తుంది'' అని చదివితే,
$\Gamma \Entails !A$ అయినప్పుడల్లా $\Gamma \Proves !A$ అని
చూపవచ్చు. నిర్దిష్ట ధర్మసంగ్రహ స్వీకృతాలను \tetoken{వ్యుత్పత్తి}{!!{derivation}} వ్యవస్థలో
చేర్చాలంటే, ఆ స్వీకృతాలను సంతృప్తిపరచే ద్వితీయ-స్థాయి నిర్మాణాలకే
అర్థవిచారాన్ని పరిమితం చేయాలి. ఉదాహరణకు, $\Delta$ ఒక ధర్మసంగ్రహ
స్వీకృతాల సమితి (బహుశా వాటన్నింటి సమితి) అయితే,
$\Gamma \cup \Delta \Entails !A$ అయినప్పుడు
$\Gamma \cup \Delta \Proves !A$ అవుతుంది. ముఖ్యంగా, మనం పరిగణిస్తున్న
ధర్మసంగ్రహ స్వీకృతాలను వాడి $!A$ను నిరూపించలేకపోతే, ఆ స్వీకృతాలు
సత్యమైనప్పటికీ $\lnot !A$కు నమూనా అయ్యే నిర్మాణం ఒకటి ఉంటుంది.

బలహీన అర్థవిచారానికి సంపూర్ణతా సిద్ధాంతం వర్తిస్తుందని చూడటానికి
అత్యంత సులభమైన మార్గం, ద్వితీయ-స్థాయి తర్కాన్ని కింది విధంగా బహు-రక
తర్కంగా భావించడం. ఒక రకానికి సాధారణ ``మొదటిస్థాయి'' \tetoken{వ్యక్తి క్షేత్రాన్ని}{!!{domain}}
అర్థనిర్దేశం చేస్తాం; తరువాత ప్రతి $k$కు ``అరిటీ $k$ గల సంబంధాల''
\tetoken{ఒక వ్యక్తి క్షేత్రం}{!!a{domain}} ఉంటుంది. భాషలో అంతర్నిర్మిత సంబంధ సంకేతాలు
``$\Atom{\Obj{true}_k}{R,x_1,\dots,x_k}$'' ఉన్నాయని తీసుకుంటాం. ఆ
సంకేతం $x_1$, \dots,~$x_k$లకు $R$ వర్తిస్తుందని ప్రకటిస్తుంది; ఇక్కడ
$R$ అనేది ``$k$-స్థాన సంబంధం'' అనే రకపు చరం; $x_1$, \dots,~$x_k$లు
మొదటిస్థాయి రకపు వస్తువులు.

ఇలా గుర్తించినప్పుడు, బలహీన ద్వితీయ-స్థాయి అర్థవిచారం సారాంశంలో
బహు-రక తర్కపు సాధారణ అర్థవిచారమే. బహు-రక తర్కాన్ని మొదటిస్థాయి
తర్కంలో ఇమిడ్చవచ్చని ఇప్పటికే గమనించాం. అందువల్ల రెండు దిశల
అనువాదాలను మినహాయిస్తే, ద్వితీయ-స్థాయి తర్కపు బలహీన భావన నిజానికి
మొదటిస్థాయి తర్కమే మారువేషంలో ఉన్న రూపం; ఇందులో తగిన స్వీకృతాలు
నియంత్రించే ``వస్తువులు'', ``సంబంధాలు'' రెండూ \tetoken{వ్యక్తి క్షేత్రంలో}{!!{domain}} ఉంటాయి.

\end{document}
